\chapter{Datenbasis und Methodik}
\label{chap:methodik}

\section{Datenbeschaffung und Aufbereitung}

Die Rohdaten wurden über die speedrun.com API~v1~\cite{speedruncom_api} abgerufen. Ausgewählt wurden 17~Spiele aus sieben~Genres (Platformer, Action-Adventure, RPG, FPS, Puzzle, Sandbox, Arcade), die eine breite stilistische und zeitliche Variation abdecken (Tabelle~\ref{tab:spiele}). Alle Analysen beschränken sich auf die \emph{any\%}-Kategorie,\footnote{Any\% bezeichnet in der Speedrunning-Community den schnellsten Weg zum Spielabschluss, unabhängig vom Fertigstellungsgrad.} da diese als einzige Kategorie in allen 17~Spielen vorhanden und zwischen Spielen vergleichbar ist.

\begin{table}[htb]
  \centering
  \caption{Spieleliste nach Genre}
  \label{tab:spiele}
  \small
  \begin{tabular}{ll}
    \toprule
    \textbf{Genre} & \textbf{Spiele} \\
    \midrule
    Platformer       & Super Mario Bros., Super Mario~64, Celeste \\
    Action-Adventure & Super Metroid, Zelda: Ocarina of Time, Hollow Knight \\
    RPG              & Pokémon Red/Blue, Final Fantasy~VII \\
    FPS              & Doom~3, Quake, Half-Life~2 \\
    Puzzle           & Portal, Portal~2, The Talos Principle \\
    Sandbox          & Minecraft: Java Edition \\
    Arcade           & Pac-Man, Donkey Konga \\
    \bottomrule
  \end{tabular}
\end{table}

Die API liefert für jedes Spiel alle verifizierten Runs der gewählten Kategorie als JSON. Für Spiele mit mehr als 10.000~Runs (Minecraft) wurde die Paginierung in umgekehrter Reihenfolge wiederholt, um vollständige Daten zu erhalten. Im Bereinigungsschritt wurden unplausible Einträge entfernt (z.\,B.\ Laufzeiten von $0$~Sekunden oder Runs ohne Zeitstempel). Die bereinigten Daten umfassen fünf~CSV-Dateien mit insgesamt 842~Weltrekord-Datenpunkten für die Lebensdaueranalyse (Q3).

\section{Statistische Methoden}

\subsection*{F1 — Prozentuale Zeitreduktion (Q1)}

Die prozentuale Zeitreduktion wird definiert als:
\[
  r_i = \frac{t_{\text{first},i} - t_{\text{best},i}}{t_{\text{first},i}} \cdot 100\,\%
\]
wobei $t_{\text{first},i}$ die erste und $t_{\text{best},i}$ die aktuell beste Laufzeit im Datensatz für Spiel~$i$ ist. Zusätzlich werden die jährliche Verbesserungsrate und die Weltrekord-Dichte berechnet.

Um zu prüfen, ob Genres sich systematisch in ihren jährlichen Raten unterscheiden, wird der \textbf{Kruskal-Wallis-Test}~\cite{kruskal1952} eingesetzt. Dieser nichtparametrische Test ist eine rangskalierte Erweiterung der einfaktoriellen Varianzanalyse und setzt keine Normalverteilung der Daten voraus --- eine wichtige Eigenschaft, da Laufzeit-Daten typischerweise stark rechtsschief verteilt sind.

\subsection*{F2 — Sättigung der Verbesserungsrate (Q2)}

Für jedes Spiel wird die kumulierte Bestzeit gegen den Weltrekord-Index aufgetragen. Anschließend werden fünf Kurvenmodelle mittels nichtlinearer Regression angepasst: logarithmisch, quadratisches Polynom~(poly2), exponentieller Zerfall, Gompertz-Kurve und LOWESS (nichtparametrisch). Der Modellvergleich erfolgt über das \textbf{Akaike-Informationskriterium}~(AIC)~\cite{akaike1974}: Ein niedrigeres AIC zeigt ein besseres Verhältnis von Anpassungsgüte zu Modellkomplexität an — komplexere Modelle werden also nur bevorzugt, wenn sie die Daten deutlich besser erklären. Das Modell mit dem niedrigsten AIC je Spiel gilt als bestes Modell.

Um den Sättigungspunkt direkt nachzuweisen, wird der \textbf{Chow-Test}~\cite{chow1960} eingesetzt. Dieser prüft, ob sich die Regressionskoeffizienten vor und nach einem möglichen Strukturbruch signifikant unterscheiden ($F$-Statistik gegen $F$-Verteilung). Alle validen Trennpunkte werden durchsucht; der Trennpunkt mit der höchsten $F$-Statistik wird als Strukturbruch gemeldet.

\subsection*{F3 — Lebensdauer von Weltrekorden (Q3)}

Die Lebensdauer eines Weltrekords wird als die Anzahl Tage gemessen, die vergehen, bis er durch einen neuen Rekord abgelöst wird. Für Rekorde, die zum Analysezeitpunkt noch Bestand haben, liegt eine \emph{rechtszensierte} Beobachtung vor: Die tatsächliche Lebensdauer ist unbekannt, aber mindestens so lang wie der bisher vergangene Zeitraum.

Der \textbf{Kaplan-Meier-Schätzer}~\cite{kaplan1958} berücksichtigt diese zensierten Beobachtungen korrekt und liefert für jeden Zeitpunkt~$t$ die geschätzte Wahrscheinlichkeit, dass ein Weltrekord bis dahin noch Bestand hat. Der Median dieser Überlebensfunktion gibt an, nach wie vielen Tagen die Hälfte aller Rekorde eines Genres bereits gebrochen wurde.

Zum Vergleich der Lebensdauerverteilungen zwischen Genres wird der \textbf{Kruskal-Wallis-Test} eingesetzt. Als Effektgröße wird $\eta^2 = (H - k + 1)/(N - k)$ berechnet~\cite{cohen1988}. Paarweise Unterschiede werden mit dem \textbf{Mann-Whitney-U-Test}~\cite{mann1947} geprüft. Die Konzentration der Verbesserungsgrößen wird über den \textbf{Gini-Koeffizienten} quantifiziert ($G = 0$: vollkommen gleiche Verbesserungen, $G = 1$: eine einzelne Verbesserung dominiert alles).
